\section{Introducción}

Un estudiante universitario reparte su semestre entre varias materias con cargas distintas, un horario de clases fijo y entregas (talleres, laboratorios, proyectos y parciales) que aparecen a lo largo del periodo. Casi siempre esa planeación se hace a mano, con agendas o calendarios genéricos que no saben cuántas horas exige realmente cada materia ni qué tan difícil es cada entrega. Este documento presenta SemestrIA, un asistente que convierte esa información en un plan semanal explicable, y recoge los tres hitos del proyecto: el problema y su contexto (Hito 1), el diseño y el avance de la solución (Hito 2) y la solución final con su evaluación (Hito 3).

\subsection{Tres problemas que nos interesaba resolver}

Antes de elegir, el equipo consideró tres problemas del entorno académico en los que la inteligencia artificial podría ayudar a un estudiante:

\begin{enumerate}
    \item \textbf{Planeación del tiempo de estudio.} Distribuir las horas de trabajo de una semana, y de las siguientes, entre materias y entregas, respetando el horario de clases, la disponibilidad personal y la dificultad de cada trabajo.
    \item \textbf{Alerta temprana de riesgo académico.} Detectar, a partir de señales de desempeño (notas parciales, entregas atrasadas, asistencia), a los estudiantes con probabilidad alta de reprobar una materia para ofrecerles apoyo a tiempo.
    \item \textbf{Orientación en la carga de materias.} Recomendar cuántas y cuáles materias inscribir en un semestre según el historial académico y la carga de trabajo que implican.
\end{enumerate}

\subsection{¿Por qué como proyecto de PTIA?}

Los tres problemas encajan en el curso porque no se resuelven con un cálculo fijo: hay restricciones que se cruzan, incertidumbre sobre el comportamiento del estudiante y decisiones que deben poder explicarse. En particular, la planeación del tiempo, que es el problema que finalmente se eligió (sección~\ref{sec:seleccion}), reúne los temas centrales de la asignatura:

\begin{itemize}
    \item \textbf{Agentes y planificación:} se trata de buscar un plan (qué hacer y cuándo) que cumpla restricciones y optimice un criterio, un problema clásico de planificación y programación de tareas \parencite{Russell2021}.
    \item \textbf{Conocimiento y razonamiento:} la equivalencia entre créditos y horas, la paridad de los laboratorios o los márgenes según la dificultad son conocimiento del dominio que el sistema usa para razonar y para justificar cada recomendación.
    \item \textbf{Aprendizaje a partir de ejemplos:} las horas que una persona tarda en un trabajo no son iguales para todos; el sistema puede aprenderlas de ejemplos que el propio estudiante registra.
    \item \textbf{Ética y responsabilidad:} un asistente que organiza el tiempo de personas afecta su bienestar y maneja datos académicos; esto obliga a pensar en autonomía, privacidad, transparencia y carga saludable.
\end{itemize}

Los problemas 2 y 3 requieren datos históricos de muchos estudiantes, que son sensibles y a los que el equipo no tiene acceso. El problema 1, en cambio, puede construirse, probarse y evaluarse con la información que el propio equipo posee (su horario y sus materias), lo que permite una prueba de concepto funcional dentro del semestre.

\subsection{Organización del documento}

La sección~\ref{sec:trabajos} revisa trabajos relacionados con cada problema; la sección~\ref{sec:seleccion} explica los criterios de selección y el problema elegido; la sección~\ref{sec:just} lo justifica; la sección~\ref{sec:alcance} delimita el alcance y los objetivos; la sección~\ref{sec:diseno} describe el diseño y la arquitectura; la sección~\ref{sec:resultados} analiza los resultados de la evaluación y la sección~\ref{sec:conclusiones} presenta las conclusiones.
